%LTeX: language=fr-FR
\documentclass[a4paper,11pt]{article}

% ---------- Préambule minimal ----------
\usepackage[french]{babel}
\usepackage{amsmath,amssymb}
\usepackage{enumitem}
\usepackage{tikz}
\usepackage{tkz-tab}
\usepackage{tasks}
\usetikzlibrary{arrows.meta}
\usepackage[margin=2.5cm]{geometry}
\usepackage{needspace}
\usepackage{currfile}
\usepackage{fancyhdr}
\usepackage{lastpage}
\settasks{label=\arabic*.}

\pagestyle{empty}
\graphicspath{ {../Images/} }
\input{\currfiledir ../def.tex}
\def \Document {Exercices}


% Environnement "exercice" : \begin{exercice}{<numéro>} ... \end{exercice}
\newenvironment{exercice}[1]{%
  \Needspace{6\baselineskip}\par\bigskip\noindent\textbf{Exercice #1}\par\smallskip\noindent\ignorespaces
}{%
  \par
}

% Repère quadrillé : \repere{xmin}{xmax}{ymin}{ymax}
\newcommand{\repere}[4]{%
  \draw[very thin, gray!40] (#1,#3) grid (#2,#4);
  \draw[->] (#1,0) -- ({#2+0.4},0) node[below right] {\small $x$};
  \draw[->] (0,#3) -- (0,{#4+0.4}) node[above left] {\small $y$};
  \foreach \x in {#1,...,#2} {\ifnum\x=0\else\draw (\x,0.08) -- (\x,-0.08) node[below] {\scriptsize $\x$};\fi}
  \foreach \y in {#3,...,#4} {\ifnum\y=0\else\draw (0.08,\y) -- (-0.08,\y) node[left] {\scriptsize $\y$};\fi}
  \node[below left] at (0,0) {\scriptsize $0$};
}

% Tableau de signes vide : \tabsignes{f}
\newcommand{\tabsignes}[1]{%
  \begin{tikzpicture}
    \tkzTabInit[lgt=1.2,espcl=2.2]{$x$/0.8,$#1(x)$/0.8}{$-\infty$,$+\infty$}
  \end{tikzpicture}%
}

% Espace vertical pour la rédaction
\newcommand{\espace}{\rule[-2cm]{0pt}{0pt}}

\begin{document}
\pagestyle{fancy}

\fancyhf{} % clear existing header/footer entries
\fancyhead[L]{ {\Niveau}}
\fancyhead[C]{ {\Chapitre}}
\fancyhead[R]{ {\Document}}
\fancyfoot[L]{ A. Fernandez }
\fancyfoot[R]{{\thepage\  / \pageref{LastPage}}}

% =====================================================
\begin{exercice}{1}
Pour chaque fonction, dire si elle est affine. Si oui, donner $a$ et $b$, et préciser si elle est linéaire ou constante.
\begin{tasks}(2)
\task $f(x)=4x+1$
\task $g(x)=-3x$
\task $h(x)=x^2-2$
\task $k(x)=5-x$
\task $m(x)=6$
\task $n(x)=2(x+3)$
\end{tasks}
\end{exercice}

% =====================================================
\begin{exercice}{2}
Soit $f$ la fonction affine définie par $f(x)=-2x+5$.
\begin{enumerate}[label=\arabic*.]
\item Calculer $f(0)$, $f(3)$, $f(-1)$ et $f(1{,}5)$.
\item Déterminer l'antécédent de $1$ par $f$, c'est-à-dire résoudre l'équation $f(x)=1$.
\item Le point $A(4~;~-3)$ appartient-il à la droite représentant $f$ ? Et le point $B(-2~;~8)$ ?
\end{enumerate}
\end{exercice}

% =====================================================
\begin{exercice}{3}
\begin{minipage}[c]{0.5\linewidth}
Dans le repère ci-contre, tracer les droites représentant les fonctions suivantes.
\begin{enumerate}[label=\arabic*.]
\item $f_1(x)=2x-3$
\item $f_2(x)=-0{,}5x+2$
\item $f_3(x)=3$
\end{enumerate}
\end{minipage}
\hfill
\begin{minipage}[c]{0.46\linewidth}
\centering
\begin{tikzpicture}[scale=0.5]
  \repere{-5}{6}{-4}{5}
\end{tikzpicture}
\end{minipage}
\end{exercice}

% =====================================================
\Needspace{14\baselineskip}
\begin{exercice}{4}
\begin{minipage}[c]{0.5\linewidth}
Lire graphiquement l'équation réduite de chacune des droites $d_1$, $d_2$, $d_3$ et $d_4$ représentées ci-contre.

\bigskip
$d_1$ : $y=\dots\dots\dots\dots$

\bigskip
$d_2$ : $y=\dots\dots\dots\dots$

\bigskip
$d_3$ : $y=\dots\dots\dots\dots$

\bigskip
$d_4$ : $y=\dots\dots\dots\dots$
\end{minipage}
\hfill
\begin{minipage}[c]{0.46\linewidth}
\centering
\begin{tikzpicture}[scale=0.5]
  \repere{-5}{5}{-5}{6}
  \draw[thick, domain=-3:2.5] plot (\x,{2*\x+1}) node[right] {$d_1$};
  \draw[thick, domain=-3:5] plot (\x,{-\x+3});
  \node[above right] at (-3,6) {$d_2$};
  \draw[thick, domain=-5:5] plot (\x,{0.5*\x-2}) node[above] {$d_3$};
  \draw[thick] (-5,-4) -- (5,-4) node[above left] {$d_4$};
\end{tikzpicture}
\end{minipage}
\end{exercice}

% =====================================================
\begin{exercice}{5}
Dans chaque cas, calculer le coefficient directeur de la droite $(AB)$.
\begin{tasks}(2)
\task $A(1~;~2)$ et $B(3~;~8)$
\task $A(0~;~5)$ et $B(2~;~1)$
\end{tasks}
\end{exercice}

% =====================================================
\begin{exercice}{6}
Dans chaque cas, déterminer l'équation réduite de la droite $(AB)$.
\begin{tasks}(3)
\task $A(1~;~5)$ et $B(3~;~9)$
\task $A(-2~;~7)$ et $B(2~;~-1)$
\task $A(0~;~-4)$ et $B(4~;~0)$
\end{tasks}
\end{exercice}

% =====================================================
\begin{exercice}{7}
Pour chaque fonction, donner son sens de variation sur $\mathbb{R}$ en justifiant.
\begin{tasks}(4)
\task $f(x)=5x-2$
\task $g(x)=7-2x$
\task $h(x)=-4$
\task $k(x)=-\dfrac{1}{2}x+1$
\end{tasks}
\end{exercice}

% =====================================================
\begin{exercice}{8}
Dresser le tableau de signes de chaque fonction sur $\mathbb{R}$.
\begin{tasks}(4)
\task $f(x)=x-4$
\task $g(x)=2x+6$
\task $h(x)=-3x+9$
\task $k(x)=4-2x$
\end{tasks}
\end{exercice}

% =====================================================
\begin{exercice}{9}
Un taxi applique un tarif qui est une fonction affine $P$ de la distance parcourue $x$, en km.
Une course de $4$~km coûte $12$~€ et une course de $10$~km coûte $24$~€.
\begin{tasks}(2)
\task Déterminer l'expression de $P(x)$.
\task Interpréter $a$ et $b$ dans le contexte.
\task Combien coûte une course de $15$~km ?
\task Quelle distance parcourir avec $30$~€ ?
\end{tasks}
\end{exercice}

% =====================================================
\begin{exercice}{10}
La batterie d'un téléphone est chargée à $100\,\%$. Pendant une vidéo, elle perd $8\,\%$ de charge par heure.
On note $B(t)$ le niveau de charge, en \%, après $t$ heures de vidéo.
\begin{enumerate}[label=\arabic*.]
\item Justifier que $B(t)=-8t+100$.
\item Quel est le sens de variation de $B$ ? Était-ce prévisible ?
\item Au bout de combien de temps la batterie atteint-elle $20\,\%$ ?
\item Résoudre $B(t)=0$ et interpréter le résultat.
\end{enumerate}
\end{exercice}

% =====================================================
\begin{exercice}{11}
Une salle d'escalade propose deux formules : la \textbf{formule A} à $12$~€ l'entrée, et la \textbf{formule B} avec une carte annuelle à $60$~€, puis $7$~€ l'entrée.
On note $x$ le nombre d'entrées dans l'année, $f(x)$ le prix payé avec la formule A et $g(x)$ le prix payé avec la formule B.

\medskip
\noindent
\begin{minipage}[c]{0.52\linewidth}
\begin{enumerate}[label=\arabic*.]
\item Exprimer $f(x)$ et $g(x)$ en fonction de $x$. Quelle est la nature de ces fonctions ?
\item Tracer les droites représentant $f$ et $g$ dans le repère ci-contre.
\item Lire graphiquement le nombre d'entrées pour lequel les deux formules coûtent le même prix.
\item Retrouver ce résultat par le calcul, en résolvant $f(x)=g(x)$.
\item À partir de combien d'entrées la formule B est-elle la plus avantageuse ?
\item Léa dispose de $200$~€. Avec quelle formule peut-elle faire le plus d'entrées ?
\end{enumerate}
\end{minipage}
\hfill
\begin{minipage}[c]{0.44\linewidth}
\centering
\begin{tikzpicture}[x=0.25cm, y=0.017cm]
  \draw[very thin, gray!40, xstep=1, ystep=20] (0,0) grid (24,300);
  \draw[thin, gray!70, xstep=4, ystep=100] (0,0) grid (24,300);
  \draw[->] (0,0) -- (25.5,0);
  \node[below=14pt] at (12,0) {\small nombre d'entrées};
  \draw[->] (0,0) -- (0,315) node[above] {\small prix (€)};
  \foreach \x in {4,8,...,24} \node[below] at (\x,0) {\scriptsize $\x$};
  \foreach \y in {100,200,300} \node[left] at (0,\y) {\scriptsize $\y$};
  \node[below left] at (0,0) {\scriptsize $0$};
\end{tikzpicture}
\end{minipage}
\end{exercice}

\end{document}
